\chapter{Multi-Agent Systems}
\label{ch:multi-agent}

\begin{quote}\itshape
Every supervisor turn is a full model call. Multi-agent is a cost decision before
it is an architecture decision.
\end{quote}

\noindent
Multi-agent is the most over-recommended pattern in the field. It is also
genuinely the right answer in a narrow set of cases, and this chapter is about
telling the two apart with numbers rather than taste.

The cost question splits cleanly in two, and keeping them apart is most of the
clarity available here.

\textbf{Structural cost} --- how many model calls a topology makes, and how much
context each carries --- is a property of the graph, not of the model. It is
deterministic, and \S\ref{sec:structural-cost} measures it exactly, for nothing,
with a counting fake.

\textbf{Quality} --- whether a supervisor actually answers better --- needs a
real model on a real task. That experiment is specified in
\S\ref{sec:quality-experiment} and \emph{has not been run}. Every claim in this
chapter that depends on it is labelled as judgement.

% ============================================================
\section{When, and the four symptoms}
\label{sec:when-multi-agent}
\index{multi-agent!when to use}

You reach for multi-agent on symptoms, not by default:

\begin{itemize}[leftmargin=1.4em]
\item \textbf{One agent's prompt has outgrown control.} The system prompt has
      become a document with sections, and changing one section breaks another.
\item \textbf{Too many tools for reliable selection.} Chapter~\ref{ch:tools}
      puts the degradation somewhere around fifteen to twenty; partitioning by
      agent is one of the three structural answers.
\item \textbf{You need genuine parallelism} over independent sub-tasks, and
      \api{Send} within one agent is not enough because the branches need
      different instructions.
\item \textbf{You need an independent critic.} A reviewer that shares the
      author's context is not independent, and that is a real limitation rather
      than a stylistic one.
\end{itemize}

Absent one of those, a pipeline with a validator is cheaper, faster and
enormously easier to debug. The default is not multi-agent; the default is one
agent, and the burden of proof is on the topology.

% ============================================================
\section{The five topologies}
\label{sec:topologies}
\index{multi-agent!topologies}
\index{supervisor}
\index{swarm}
\index{orchestrator}
\index{pipeline}

\mermaidfig{ma-topologies}{%
  All five to the same scale, shaded by structural cost: green cheapest, amber
  dear, red hardest to reason about.}{%
  The five multi-agent topologies compared.}

\begin{center}
\small
\begin{tabularx}{\linewidth}{@{}lXX@{}}
\toprule
\textbf{Topology} & \textbf{Shape} & \textbf{Reach for it when} \\
\midrule
Pipeline & Fixed sequence, no routing call & The steps are known in advance.
  Cheapest by a distance. \\[3pt]
Swarm & Agents hand control directly to each other & A conversation moves
  between specialisms. \\[3pt]
Supervisor & A router chooses a worker, collects, repeats & You need control and
  an audit trail of who did what. \\[3pt]
Orchestrator-worker & Planner splits, workers run in parallel, a synthesiser
  merges & Research, or N documents to process. \\[3pt]
Network & Everyone may call everyone & Rarely. Hardest to trace, easiest to
  loop. \\
\bottomrule
\end{tabularx}
\end{center}

\begin{note}
Note what a pipeline does \emph{not} have: a routing model call. The sequence is
in your code. That is why it is the cheapest topology and why it should be the
first thing you try --- most problems described as needing a supervisor need a
pipeline and a validator.
\end{note}

% ============================================================
\section{Handoffs}
\label{sec:handoffs}
\index{handoff}
\index{Command@\texttt{Command}!PARENT}

A handoff is Chapter~\ref{ch:stategraph}'s \api{Command}: update state and jump,
in one return. Between siblings in one graph it is an ordinary \code{goto}.
Across a subgraph boundary it needs \code{Command.PARENT}:

\begin{python}[caption={An agent inside a subgraph handing control to a sibling of its parent}, label={lst:handoff}]
def billing_agent(state) -> Command:
    return Command(
        goto="technical",
        graph=Command.PARENT,
        update={"trail": ["billing handed off"]},
    )
\end{python}

\begin{console}
{'trail': ['billing handed off', 'technical handled it']}
\end{console}

That is the swarm pattern in one construct: the billing agent decided, on its
own, that the technical agent should take over, and control moved without a
router being consulted.

\begin{warning}
Remember Chapter~\ref{ch:stategraph}'s finding: annotate the return type as
\code{Command[Literal["technical", ...]]} or the handoff edge does not appear in
the drawn graph. In a multi-agent system that matters more than anywhere else,
because the diagram is the only practical way to see the topology --- and an
unannotated swarm draws as a set of disconnected agents.
\end{warning}

\subsection{Handoff loops}
\label{sec:handoff-loops}
\index{handoff!loops}

The failure mode unique to swarms:

\mermaidfig{ma-handoff-loop}{%
  Two agents, each correctly identifying the other half of a mixed question as
  outside its remit.}{%
  A handoff loop, and why neither agent is wrong.}

Neither agent is malfunctioning. A user who says ``I was charged twice and the
app crashes'' has asked a question that genuinely spans both, and each agent
correctly identifies the other half as somebody else's. The loop is a property of
the routing rule, not of the model.

Three defences, in order of preference: give one agent explicit ownership of
mixed questions; count hops in state and force a decision after N; and set
Chapter~\ref{ch:persistence}'s \code{recursion\_limit} as the backstop it is,
rather than as a fix.

% ============================================================
\section{Structural cost, measured}
\label{sec:structural-cost}
\index{multi-agent!token economics}
\index{measurement!structural cost}

Here is the part that does not need a provider.

\textbf{Method.} One fixed question. Five implementations, each answering it. The
model is a fake that records every call --- how many, and how many characters of
context each carried --- before replying from a fixed script. Because the script
is fixed, the numbers are exactly reproducible; because call count is a property
of the graph, a real model would make the same number of calls. The script is
\code{code/ch11/structural\_cost.py} and runs in under a second.

\begin{center}
\small
\begin{tabularx}{\linewidth}{@{}lrrrX@{}}
\toprule
\textbf{Topology} & \textbf{Model calls} & \textbf{Context chars} & \textbf{vs baseline} \\
\midrule
Single agent & 2 & 172 & 1.0$\times$ \\
Pipeline (3 stages) & 3 & 243 & 1.4$\times$ \\
Swarm (3 hops) & 3 & 240 & 1.4$\times$ \\
Orchestrator + 3 workers & 5 & 393 & 2.3$\times$ \\
Supervisor (2 workers) & 5 & 487 & 2.8$\times$ \\
\bottomrule
\end{tabularx}
\end{center}

\begin{warning}
Read these two columns with different levels of trust.

\textbf{Model calls is a hard number.} It is determined by the topology, and a
real provider makes exactly the same count.

\textbf{Context chars is indicative.} It depends on how long the scripted replies
are, and real replies are longer and vary. The \emph{ratios} hold for the
structural reason given below; the absolute figures are an artefact of the
script and should not be quoted.
\end{warning}

Three things are worth drawing out.

\textbf{Supervisor costs more than orchestrator despite the same call count.}
Both make five calls; the supervisor carries 2.8$\times$ the context against
2.3$\times$. The reason is structural: every router turn re-reads the whole
accumulating conversation, so a supervisor's context grows with each delegation
while an orchestrator's workers each see only their own slice.

\mermaidfig{ma-supervisor-turns}{%
  Where a supervisor's five calls go. Three of them are routing rather than
  work.}{%
  Model calls per user turn under a supervisor.}

\textbf{Swarm and pipeline cost essentially the same} --- 240 against 243. A
handoff is not structurally more expensive than a fixed edge. What you pay for
with a swarm is not tokens; it is traceability.

\textbf{The single agent is the baseline for a reason.} Two calls. Every topology
in the table is a multiple of it, and the multiple is the thing you have to
justify.

\begin{note}
A useful way to hold this: with a supervisor, \textbf{three of five model calls
are routing rather than work}. You are paying a full model call to decide who
should do something, twice, plus one to compose the result. If the routing
decision is obvious from the input --- and it very often is --- that is three
calls spent on a problem a conditional edge solves for free.
\end{note}

\subsection{Model selection per role}
\label{sec:model-per-role}
\index{cost!model per role}

The structural numbers suggest the obvious mitigation. Routing is a
classification task and rarely needs the strongest model available; the work is
where capability matters.

Putting a cheap model on the supervisor and a strong one on the workers inverts
the usual advice, and the measurement is why: the supervisor makes the most calls
and carries the most context, so it is the most expensive seat in the topology.

\begin{warning}
That is a hypothesis with a cost model behind it, not a measured result. Whether
a cheap router degrades routing quality enough to cost more in retries than it
saves per call is exactly the sort of question \S\ref{sec:quality-experiment}
exists to answer, and it has not been run. Treat the recommendation as judgement.
\end{warning}

% ============================================================
\section{The quality experiment}
\label{sec:quality-experiment}
\index{measurement!multi-agent benchmark}

\begin{warning}
\textbf{This experiment has not been run.} Everything below is the specification,
not results. The tables in Appendix~\ref{app:cheatsheet} are deliberately empty
and stay empty until it has been. No comparative quality claim anywhere in this
chapter is stated as fact.
\end{warning}

\textbf{Design.} One fixed task with a checkable answer. Six implementations: the
single agent baseline and the five topologies. Model, temperature, instructions
and tool set held constant across all six --- only the topology varies.

\textbf{Metrics.} Turn count; token cost from \code{usage\_metadata}
(Chapter~\ref{ch:models}); wall-clock latency; and success rate against the
checkable answer.

\textbf{Discipline.} At least twenty runs per implementation. Report medians and
spreads, never best runs --- a topology that occasionally produces a brilliant
answer and usually produces a poor one is worse than a consistent one, and a best
run hides exactly that.

\begin{note}
\textbf{Keep the raw event streams, not just the summary rows.}
Chapter~\ref{ch:observability} re-scores these same runs under an evaluation
harness --- trajectory quality, tool-selection accuracy, whether the critic
contributed anything --- and re-running to recover traces is expensive. Store
them at the point of collection.

The structural numbers in \S\ref{sec:structural-cost} give this experiment a
prediction to test: if the supervisor does not score materially better than the
single agent, it is spending 2.5 times the calls for nothing, and the chapter's
recommendation becomes concrete rather than cautious.
\end{note}

% ============================================================
\section{Debugging a multi-agent run}
\label{sec:debugging-multi-agent}
\index{debugging!multi-agent}
\index{troubleshooting!namespace}

The honest difficulty: when a five-agent system gives a poor answer, the question
``which agent went wrong'' is genuinely hard, because each agent's output is the
next one's input and a small early error compounds.

Three things make it tractable, all of them from earlier chapters.

\textbf{Stream with subgraph namespaces.} Chapter~\ref{ch:stategraph}'s
\code{subgraphs=True} tells you which level produced each update. Without it a
composed graph reports only aggregate updates from its top-level nodes.

\textbf{Name your agents.} \api{create\_agent} takes a \code{name}, and it
appears in traces and in the \code{subagents} stream projection
(Chapter~\ref{ch:interrupts}). Unnamed agents in a trace are indistinguishable.

\textbf{Keep a trail in state.} An \code{Annotated[list, operator.add]} field
that each agent appends its name to costs nothing and answers ``what actually
ran, in what order'' without a tracing backend. Every listing in this chapter
uses one.

% ============================================================
\section{The functional API}
\label{sec:functional-api}
\index{entrypoint@\texttt{@entrypoint}}
\index{task@\texttt{@task}}

Not every workflow wants to be a graph. When the control flow is naturally
sequential and expressible in ordinary Python, \api{@entrypoint} and \api{@task}
give you the same runtime without the builder:

\begin{python}[caption={Same runtime, ordinary control flow}, label={lst:functional}]
from langgraph.func import entrypoint, task


@task
def analyse(doc: str) -> str:
    return expensive_analysis(doc)


@entrypoint(checkpointer=checkpointer)
def workflow(docs: list[str]) -> list[str]:
    futures = [analyse(d) for d in docs]      # scheduled
    return [f.result() for f in futures]      # resolved
\end{python}

\begin{console}
['A', 'B', 'C']
type of workflow: Pregel
state after run : ['A', 'B', 'C']
\end{console}

The object is a \api{Pregel} --- the same runtime as a compiled graph, verified by
\code{get\_state} returning the run's state afterwards. So checkpointing,
interrupts and streaming all apply. \api{@task} takes the same
\code{retry\_policy}, \code{cache\_policy} and \code{timeout} as a node
(Chapter~\ref{ch:persistence}).

\begin{note}
The trade is visibility. A graph can be drawn, validated for reachability and
inspected before it runs; a function cannot. Choose the functional API when the
flow is genuinely linear and the branching is ordinary Python that would be
awkward as edges --- and choose the graph when the topology is the thing you want
to see, which in a multi-agent system it usually is.
\end{note}

% ============================================================
\section{Seeing the topology}
\label{sec:visualisation}
\index{workflow!visualisation}

Print the graph. It is one line, it is free, and by this chapter the answer stops
being obvious:

\begin{python}[caption={The topology, from the running code}, label={lst:draw}]
print(graph.get_graph().draw_mermaid())
print(graph.get_graph(xray=True).draw_mermaid())   # expand subgraphs too
\end{python}

\code{langgraph dev} serves a local development UI over a \code{langgraph.json},
with LangGraph Studio for stepping through a run and inspecting state at each
node. It is the most useful thing in the ecosystem for understanding somebody
else's graph, and it is worth ten minutes.

\needscreenshot{studio-multi-agent}{%
  LangGraph Studio showing a supervisor topology mid-run.}{%
  Run a supervisor graph under langgraph dev and open Studio. Capture the graph
  panel with the supervisor node highlighted as the active one, with at least one
  worker already completed, and the state inspector open on the right showing the
  message list. The point of the capture is that the topology and the current
  position are visible together, which is what a printed diagram cannot show.}

% ============================================================
\section{What to take away}
\label{sec:ch11-summary}

\begin{itemize}[leftmargin=1.4em]
\item Multi-agent is a response to symptoms --- an unmanageable prompt, too many
      tools, real parallelism, an independent critic. Not a default.
\item Five topologies. A pipeline has no routing call and is the cheapest thing
      that could work; try it first.
\item Handoffs are \api{Command}, and across a subgraph boundary
      \code{Command.PARENT}. Annotate the return type or the topology does not
      draw.
\item Handoff loops are a routing-rule defect, not a model defect. Give one agent
      ownership of mixed questions.
\item \textbf{Measured structural cost}, single agent as baseline: pipeline and
      swarm 3 calls, orchestrator and supervisor 5. Supervisor carries the most
      context because every router turn re-reads the whole conversation.
\item With a supervisor, three of five model calls are routing rather than work.
\item Model calls is a hard number; context size is indicative and
      script-dependent.
\item The quality half is unrun. Every comparative quality claim here is
      judgement, and Appendix~\ref{app:cheatsheet} stays empty until it is not.
\item The functional API is the same \api{Pregel} runtime with ordinary control
      flow, trading visibility for directness.
\end{itemize}

\begin{exercisebox}
\textbf{1.} Run \code{code/ch11/structural\_cost.py}. Then add a third worker to
the supervisor and predict both columns before re-running. The call count should
be exactly predictable; check whether your context prediction was as good.

\textbf{2.} Build the handoff loop deliberately: two agents, a question that
spans both, no hop limit. Watch it run until \code{recursion\_limit}. Then add
a hop counter in state and force a decision.

\textbf{3.} Take a supervisor graph and print it with and without
\code{xray=True}. Then remove the \code{Command} return annotations and print
again. Keep all three.

\textbf{4.} Implement the same small task as a pipeline and as a supervisor.
Measure both with the harness. Then ask whether the supervisor bought anything
you could name --- and notice that you cannot answer that without
\S\ref{sec:quality-experiment}.

\textbf{5.} Rewrite one topology with \api{@entrypoint} and \api{@task}. Confirm
\code{get\_state} still works, then try to draw it.
\end{exercisebox}

\begin{projectbox}
\textbf{Stage 08 --- the supervisor variant, and the comparison.} A supervisor
build of Ops Copilot alongside the single-agent one, and the harness that runs
both on one fixed task.

Two requirements in \code{stages/08-multi-agent/README.md} carry the weight. The
harness must report \textbf{medians and spreads over at least twenty runs}, never
a best run. And it must persist \textbf{raw event streams}, not only summary
rows, because Chapter~\ref{ch:observability} re-scores these same runs and
re-running to recover traces is expensive.

The structural half of the comparison can be built and tested immediately with
the metered fake from \S\ref{sec:structural-cost} --- no provider, no budget, and
it belongs in CI. The quality half waits for the experiment.
\end{projectbox}
